$\psi$ is strictly convex, being $c\,2^{t}$ up to an affine term, and $\psi(0)=0$; a strictly
convex function has at most two roots, so at most one root is non-zero, and
$\psi(t)\rightarrow+\infty$ at both ends. The sign of $\psi'(0)=c\ln2-1$ therefore decides
where the second root lies: to the right of $0$ when $c<1/\ln2$, to the left when
$c>1/\ln2$, and at $c=1/\ln2$ the point $0$ is the strict minimum and there is no other root.
The second equivalence of \emph{(1)} is the same convexity read at $t=1$, where
$\psi(1)=c-1$: for $c<1$ the point $1$ lies strictly between the two roots, so $p_*>1$; for
$c=1$ it is itself a root; for $c>1$ it lies outside their interval, so $p_*<1$.

Item \emph{(2)} is $\psi(p_*)=0$, that is $c(\tau-1)=p_*$, divided by $p_*\neq0$; the
integral form is the same identity read as
$\int_0^1c\ln2\cdot2^{p_*v}\,dv=c(\tau-1)/p_*=1$. For \emph{(3)}, $0<c<1$ gives $p_*>1$ and
hence $\tau>2$; the strictly convex $t\mapsto2^{t}-t-1$ vanishes at $t=0$ and $t=1$ and is
therefore positive beyond $1$, and evaluating it at $p_*$, $p_*+1$ and $p_*+2$ gives the three
inequalities of \eqref{eq:doubling_cramer_ineq}.
